%% Included in decks: slidemage-manual-copy
%%

%% Translation of: neg-EN-Installing-LaTeX-Linux
%%

\begin{standardframe}{LaTeX unter Linux installieren}
	\small
	SlideMage übersetzt jedes Deck mit LuaLaTeX über latexmk und braucht dafür
	eine TeX-Distribution. Unter Linux ist das TeX Live, entweder aus der
	Paketverwaltung der Distribution oder mit dem offiziellen Installer.
	\par
	\begin{columns}[T,onlytextwidth]
		\begin{column}{0.56\linewidth}
			\shrinklists{1}
			\textbf{Paketverwaltung} (am einfachsten)
			\begin{itemize}
				\item Debian/Ubuntu: \SMCode{sudo apt install texlive-full latexmk}
				\item Fedora: \SMCode{sudo dnf install texlive-scheme-full latexmk}
				\item Arch: \SMCode{sudo pacman -S texlive texlive-lang}
				\item Oft ein, zwei Jahre hinter dem aktuellen TeX Live, für SlideMage reicht das
			\end{itemize}
		\end{column}
		\begin{column}{0.4\linewidth}
			\shrinklists{1}
			\textbf{TeX Live direkt} (\SMCode{tug.org/texlive})
			\begin{itemize}
				\item \SMCode{install-tl}, Schema \emph{full} (etwa 8\,GB), im Home-Verzeichnis ohne Root-Rechte
				\item Den \SMCode{bin}-Ordner in den \SMCode{PATH}
			\end{itemize}
			\medskip
			\textbf{In beiden Fällen}
			\begin{itemize}
				\item Im Terminal testen: \SMCode{lualatex --version}, \SMCode{latexmk -v}
			\end{itemize}
		\end{column}
	\end{columns}
\end{standardframe}
